\documentclass[11pt]{article}
\usepackage[utf8]{inputenc}
\usepackage{amsmath,amssymb,booktabs,geometry,verbatim}
\usepackage{CJKutf8}
\geometry{margin=1in}
\usepackage[hidelinks]{hyperref}
\makeatletter\let\@author\@empty\makeatother

\title{CSP-S 2023 (Senior)}
\date{2023-10-21, 14:30--18:30 (4 hours)}

\begin{document}
\maketitle

\noindent
Translated from the official Chinese problem booklet
(2023 CCF CSP-J/S second round, senior level).
All four problems: time limit 1 s, memory limit 512 MB, 100 points each.
Input is read from the file \texttt{<name>.in} and output is written to
\texttt{<name>.out}, where \texttt{<name>} is the problem code shown in
each heading.

\clearpage
\section*{Problem 1: Combination Lock
  (\begin{CJK}{UTF8}{gbsn}密码锁\end{CJK}, lock)
  \hfill{\normalsize\href{https://qoj.ac/problem/7813}{QOJ 7813}}}

\noindent
Time limit: 1 s.
Memory limit: 512 MB.

Little Y has a combination lock with five dials.
As shown in the figure, each dial carries the digits from 0 to 9.
Every dial is cyclic from 0 to 9, that is, turning 9 by one position gives
0 or 8.

\begin{center}
[Figure 1: a combination lock (photograph omitted).]
\end{center}

Because the campus is fairly safe, Little Y locks the bike as follows:
starting from the correct combination, the lock is turned randomly
exactly once; each time, either a single dial is turned by some amount, or
two adjacent dials are turned simultaneously.

When Little Y chooses to turn two adjacent dials simultaneously, the two
dials are turned by the same amount.
That is, Little Y can turn the lock from \texttt{0 0 1 1 5} into
\texttt{1 1 1 1 5}, but not into \texttt{1 2 1 1 5}.

As time went by, Little Y started to worry about the safety of this way of
locking, so Little Y recorded $n$ states of the lock after locking the
bike.
Note that none of these $n$ states is the correct combination.

To check the safety of this way of locking, Little Y wants to know how
many possible correct combinations there are such that each of them can
produce all $n$ recorded states of the lock by the locking method
described above.

\subsection*{Input}

Read the input from the file \texttt{lock.in}.

The first line contains a positive integer $n$, the number of states of
the lock after locking.

Each of the next $n$ lines contains five integers describing one state of
the lock.

\subsection*{Output}

Write the output to the file \texttt{lock.out}.

Output one line containing one integer: the number of correct
combinations that correspond to these $n$ states of the lock under the
given locking method.

\subsection*{Sample 1 input}

\begin{verbatim}
1
0 0 1 1 5
\end{verbatim}

\subsection*{Sample 1 output}

\begin{verbatim}
81
\end{verbatim}

\subsection*{Explanation of sample 1}

There are 81 possible combinations in total.
Among them, 45 are obtained by turning one dial and 36 are obtained by
turning two dials.

\subsection*{Sample 2}

See the files \texttt{lock/lock2.in} and \texttt{lock/lock2.ans} in the
contestant's directory.

\subsection*{Constraints}

For all test data: $1 \le n \le 8$.

\begin{center}
\begin{tabular}{ccc}
\toprule
Test cases & $n \le$ & Special property \\
\midrule
$1 \sim 3$ & 1 & none \\
$4 \sim 5$ & 2 & none \\
$6 \sim 8$ & 8 & A \\
$9 \sim 10$ & 8 & none \\
\bottomrule
\end{tabular}
\end{center}

Special property A: it is guaranteed that every correct combination can
produce the $n$ given states by turning only one dial.

\clearpage
\section*{Problem 2: Match-Three (\begin{CJK}{UTF8}{gbsn}消消乐\end{CJK}, game)
  \hfill{\normalsize\href{https://qoj.ac/problem/7814}{QOJ 7814}}}

\noindent
Time limit: 1 s.
Memory limit: 512 MB.

Little L is playing a low-budget version of a match-three game.
This version of the game is one-dimensional, and each move can only
eliminate two adjacent elements.

Now he has a string of length $n$ consisting only of lowercase letters.
We call a string \emph{erasable} if and only if it can be turned into the
empty string by performing some number of operations on it.

In each operation, two adjacent equal characters are deleted from the
string, and the remaining parts of the string are concatenated.

Little L wants to know how many of the non-empty contiguous substrings of
this string are erasable.

\subsection*{Input}

Read the input from the file \texttt{game.in}.

The first line contains a positive integer $n$, the length of the string.

The second line contains a string of length $n$ consisting only of
lowercase letters, the string asked about in the problem.

\subsection*{Output}

Write the output to the file \texttt{game.out}.

Output one line containing one integer, the answer to the question.

\subsection*{Sample 1 input}

\begin{verbatim}
8
accabccb
\end{verbatim}

\subsection*{Sample 1 output}

\begin{verbatim}
5
\end{verbatim}

\subsection*{Explanation of sample 1}

There are 5 erasable contiguous substrings in total, namely
\texttt{cc}, \texttt{acca}, \texttt{cc}, \texttt{bccb} and
\texttt{accabccb}.

\subsection*{Sample 2}

See the files \texttt{game/game2.in} and \texttt{game/game2.ans} in the
contestant's directory.

\subsection*{Sample 3}

See the files \texttt{game/game3.in} and \texttt{game/game3.ans} in the
contestant's directory.

\subsection*{Sample 4}

See the files \texttt{game/game4.in} and \texttt{game/game4.ans} in the
contestant's directory.

\subsection*{Constraints}

For all test data: $1 \le n \le 2 \times 10^6$, and the string asked
about consists only of lowercase letters.

\begin{center}
\begin{tabular}{ccc}
\toprule
Test cases & $n \le$ & Special property \\
\midrule
$1 \sim 5$ & 10 & none \\
$6 \sim 7$ & 800 & none \\
$8 \sim 10$ & 8000 & none \\
$11 \sim 12$ & $2 \times 10^5$ & A \\
$13 \sim 14$ & $2 \times 10^5$ & B \\
$15 \sim 17$ & $2 \times 10^5$ & none \\
$18 \sim 20$ & $2 \times 10^6$ & none \\
\bottomrule
\end{tabular}
\end{center}

Special property A: every character of the string is chosen independently
and uniformly at random from the alphabet.

Special property B: the string consists only of \texttt{a} and
\texttt{b}.

\clearpage
\section*{Problem 3: Structs (\begin{CJK}{UTF8}{gbsn}结构体\end{CJK}, struct)
  \hfill{\normalsize\href{https://qoj.ac/problem/7815}{QOJ 7815}}}

\noindent
Time limit: 1 s.
Memory limit: 512 MB.

\subsection*{Background}

In high-level languages such as C++, besides basic types such as
\texttt{int} and \texttt{float}, one can usually also define struct
types.
In this problem you need to simulate the way structs are defined in a
high-level language similar to C++, and compute the corresponding memory
usage and related information.

\subsection*{Description}

In this language there are 4 basic types: \texttt{byte}, \texttt{short},
\texttt{int} and \texttt{long}, occupying 1, 2, 4 and 8 bytes of space
respectively.

When defining a struct \emph{type}, one gives the \emph{type name} and
the \emph{members}, where for each member one gives, in order, its
\emph{type} and its \emph{name}.
The type may be a basic type or a \emph{previously defined} struct type.
Note that defining a struct \emph{type} does not define a concrete
element, i.e.\ it does not occupy memory.

When defining an \emph{element}, one gives the \emph{type} and the
\emph{name} of the element.
An element occupies memory according to the following rules:
\begin{itemize}
\item All members of an element are laid out in memory in the
  \emph{order in which they were given in the definition}; the same
  holds for members whose type is a struct.
\item To guarantee efficient memory access, the addresses occupied by an
  element must satisfy the \emph{alignment rule}: the \emph{size} of
  any type, and the \emph{starting address} of any element of that type
  in memory, must both be aligned to an \emph{integer multiple} of the
  alignment requirement of that type.
  Concretely:
  \begin{itemize}
  \item For a basic type, the alignment requirement equals the size of
    the space it occupies; for example, the type \texttt{int} must be
    aligned to 4 bytes, and similarly for the others.
  \item For a struct type, the alignment requirement equals the
    \emph{maximum} of the alignment requirements of its members; for
    example, a struct type containing an \texttt{int} and a
    \texttt{short} must be aligned to 4 bytes.
  \end{itemize}
\end{itemize}

Here is an example (written in C++ syntax):
\begin{verbatim}
struct d {
    short a;
    int b;
    short c;
};
d e;
\end{verbatim}

This code defines the struct type \texttt{d} and the element \texttt{e}.
The element \texttt{e} contains three members \texttt{e.a}, \texttt{e.b}
and \texttt{e.c}, which occupy the addresses of bytes $0 \sim 1$,
$4 \sim 7$ and $8 \sim 9$ respectively.
Since the type \texttt{d} must be aligned to 4 bytes, \texttt{e}
occupies the addresses of bytes $0 \sim 11$ and has size 12 bytes.

You need to process $n$ operations, each of which is one of the following
four kinds:
\begin{enumerate}
\item Define a struct type.
  Concretely, a positive integer $k$ and strings
  $s, t_1, n_1, \ldots, t_k, n_k$ are given, where $k$ is the number of
  members of the type, $s$ is the type name, $t_1, t_2, \ldots, t_k$ are
  the types of the members in order, and $n_1, n_2, \ldots, n_k$ are the
  names of the members in order.
  You need to output the size and the alignment requirement of this
  struct type, separated by one space.
\item Define an element.
  Concretely, strings $t$ and $n$ are given, the type and the name of the
  element respectively.
  All defined elements are laid out one after another in order, starting
  from memory address 0, and must satisfy the address alignment rule.
  You need to output the starting address of the newly defined element.
\item Access an element.
  Concretely, a string $s$ is given, denoting the element being accessed.
  As in C++ and similar languages, \texttt{.} is used to access a member
  of a struct type.
  For example, \texttt{a.b.c} means that \texttt{a} is a defined element
  of some struct type having a member named \texttt{b}, which is itself of
  some struct type having a member named \texttt{c}.
  You need to output the starting address of the \emph{innermost}
  element accessed in this way.
\item Access a memory address.
  Concretely, a non-negative integer $addr$ is given, the address being
  accessed.
  You need to determine whether some element of a \emph{basic type}
  occupies this address.
  If so, output this element in the element-access format of operation 3;
  otherwise output \texttt{ERR}.
\end{enumerate}

\subsection*{Input}

Read the input from the file \texttt{struct.in}.

Line 1: a positive integer $n$, the number of operations.

The following lines describe the operations in order; the first positive
integer $op$ on each line denotes the type of the operation:
\begin{itemize}
\item If $op = 1$, first a string $s$ and a positive integer $k$ are
  given, the type name and the number of members; then each of the next
  $k$ lines contains two strings $t_i$ and $n_i$, the type and the name
  of each member in order.
\item If $op = 2$, two strings $t$ and $n$ are given, the type and the
  name of the element.
\item If $op = 3$, a string $s$ is given, the element being accessed.
\item If $op = 4$, a non-negative integer $addr$ is given, the address
  being accessed.
\end{itemize}

\subsection*{Output}

Write the output to the file \texttt{struct.out}.

Output $n$ lines, the results of the operations in order; the required
output is as described in the problem description.

\subsection*{Sample 1 input}

\begin{verbatim}
5
1 a 2
short aa
int ab
1 b 2
a ba
long bb
2 b x
3 x.ba.ab
4 10
\end{verbatim}

\subsection*{Sample 1 output}

\begin{verbatim}
8 4
16 8
0
4
x.bb
\end{verbatim}

\subsection*{Explanation of sample 1}

In the struct type \texttt{a}, the member \texttt{aa} of type
\texttt{int} occupies the addresses of bytes $0 \sim 3$, and the member
\texttt{ab} of type \texttt{short} occupies the addresses of bytes
$4 \sim 5$.
Since its alignment requirement is 4 bytes, its size is 8 bytes.
In the same way one computes that the struct type \texttt{b} has size
16 bytes and alignment requirement 8 bytes.

(Translator's note: the sample explanation above is translated verbatim
from the original booklet, which lists the members of \texttt{a} as
\texttt{aa} of type \texttt{int} and \texttt{ab} of type \texttt{short},
whereas the sample input declares \texttt{short aa} and \texttt{int ab}.)

\subsection*{Sample 2}

See the files \texttt{struct/struct2.in} and \texttt{struct/struct2.ans}
in the contestant's directory.

\subsection*{Explanation of sample 2}

In the second operation of type 4, the accessed memory address lies
exactly in a ``hole'' left for address alignment, so no element of a
basic type occupies it.

\subsection*{Sample 3}

See the files \texttt{struct/struct3.in} and \texttt{struct/struct3.ans}
in the contestant's directory.

\subsection*{Constraints}

For all data, $1 \le n \le 100$, $1 \le k \le 100$ and
$0 \le addr \le 10^{18}$.

All struct type names, member names and element names consist of at most
10 lowercase letters, and none of them is \texttt{byte}, \texttt{short},
\texttt{int} or \texttt{long} (i.e.\ they do not coincide with the names
of basic types).

All struct type names and element names are pairwise distinct, and the
member names within one struct are pairwise distinct.
However, different structs may have members with the same name, and a
member name inside a struct may coincide with the name of a struct type or
of an element.

It is guaranteed that all operations conform to the rules and
requirements stated in the problem; i.e.\ a struct definition never uses a
non-existent type, no non-existent element or member is accessed, and so
on.

It is guaranteed that the size of every struct, and the highest memory
address occupied by any defined element, do not exceed $10^{18}$.

\begin{center}
\begin{tabular}{cc}
\toprule
Test case number & Special properties \\
\midrule
$1$ & A, D \\
$2 \sim 3$ & A \\
$4 \sim 5$ & B, D \\
$6 \sim 8$ & B \\
$9 \sim 10$ & C, D \\
$11 \sim 13$ & C \\
$14 \sim 16$ & D \\
$17 \sim 20$ & none \\
\bottomrule
\end{tabular}
\end{center}

Special property A: there is no operation of type 1.

Special property B: there is exactly one operation of type 1.

Special property C: all member types given in operations of type 1 are
basic types.

Special property D: the only basic type used is \texttt{long}.

\subsection*{Hint}

The alignment requirement and the size of a struct type are formally
defined as follows.
\begin{itemize}
\item Suppose the struct has $k$ members with sizes $s_1, \ldots, s_k$
  and alignment requirements $a_1, \ldots, a_k$.
\item Then the alignment requirement of the struct is
  $a = \max\{a_1, \ldots, a_k\}$.
\item Further suppose the \emph{address offsets} of these members in the
  layout are $o_1, \ldots, o_k$.
  Then:
  \begin{itemize}
  \item $o_1 = 0$;
  \item for $i = 2, \ldots, k$, $o_i$ is the smallest value such that
    $o_{i-1} + s_{i-1} \le o_i$ and $a_i$ divides $o_i$;
  \item the size $s$ of the struct is the smallest value such that
    $o_k + s_k \le s$ and $a$ divides $s$.
  \end{itemize}
\end{itemize}

The memory layout of defined elements is formally defined as follows.
\begin{itemize}
\item Suppose the $i$-th defined element has size $s_i$, alignment
  requirement $a_i$ and starting address $b_i$.
\item Then $b_1 = 0$, and for $2 \le i$, $b_i$ is the smallest value
  such that $b_{i-1} + s_{i-1} \le b_i$ and $a_i$ divides $b_i$.
\end{itemize}

\clearpage
\section*{Problem 4: Planting Trees (\begin{CJK}{UTF8}{gbsn}种树\end{CJK}, tree)
  \hfill{\normalsize\href{https://qoj.ac/problem/7816}{QOJ 7816}}}

\noindent
Time limit: 1 s.
Memory limit: 512 MB.

You are a forest ranger.
One day you receive a task: plant trees on the plots of a forest and tend
them until the trees grow to the specified heights.

The map of the forest has $n$ plots, and plot number 1 is connected to
the entrance of the forest.
There are $n - 1$ roads connecting these plots, so that every plot can be
reached from every other plot via the roads.
Initially there are no trees on any plot.

Your goal is to plant one tree on every plot, and to make the tree on
plot $i$ grow to a height of at least $a_i$ metres.

Each day you may choose a plot that has no tree yet and is
\emph{directly adjacent to some plot that already has a tree (i.e.\
connected to it by a single road)}, and plant a tree of height 0 metres
there.
If every plot has already been planted, you do nothing on that day.
In particular, on day 1 you can only plant on the empty plot number 1.

For each plot, starting from the day on which its tree is planted, the
tree on that plot grows by a certain height every day.
Because of differing climate and soil conditions, on day $x$ the tree on
plot $i$ grows by $\max(b_i + x * c_i, 1)$ metres.
Note that $x$ is counted from the first day of the whole task, not from
the first day after this tree was planted.

You want to know: what is the minimum number of days needed to complete
your task?

\subsection*{Input}

Read the input from the file \texttt{tree.in}.

The first line contains a positive integer $n$, the number of plots in
the forest.

Each of the next $n$ lines contains three integers $a_i, b_i, c_i$
describing one plot, with the meaning given in the problem description.

Each of the next $n - 1$ lines contains two positive integers $u_i, v_i$
denoting a road connecting plots $u_i$ and $v_i$.

\subsection*{Output}

Write the output to the file \texttt{tree.out}.

Output one line containing one positive integer, the minimum number of
days needed to complete the task.

\subsection*{Sample 1 input}

\begin{verbatim}
4
12 1 1
2 4 -1
10 3 0
7 10 -2
1 2
1 3
3 4
\end{verbatim}

\subsection*{Sample 1 output}

\begin{verbatim}
5
\end{verbatim}

\subsection*{Explanation of sample 1}

Day 1: plant a tree on plot 1; the tree on plot 1 grows to 2 metres.

Day 2: plant a tree on plot 3; the trees on plots 1, 3 grow to 5, 3
metres respectively.

Day 3: plant a tree on plot 4; the trees on plots 1, 3, 4 grow to 9, 6, 4
metres respectively.

Day 4: plant a tree on plot 2; the trees on plots 1, 2, 3, 4 grow to
14, 1, 9, 6 metres respectively.

Day 5: the trees on plots 1, 2, 3, 4 grow to 20, 2, 12, 7 metres
respectively.

\subsection*{Sample 2}

See the files \texttt{tree/tree2.in} and \texttt{tree/tree2.ans} in the
contestant's directory.

\subsection*{Sample 3}

See the files \texttt{tree/tree3.in} and \texttt{tree/tree3.ans} in the
contestant's directory.

\subsection*{Sample 4}

See the files \texttt{tree/tree4.in} and \texttt{tree/tree4.ans} in the
contestant's directory.

\subsection*{Constraints}

For all test data: $1 \le n \le 10^5$, $1 \le a_i \le 10^{18}$,
$1 \le b_i \le 10^9$, $0 \le |c_i| \le 10^9$ and $1 \le u_i, v_i \le n$.
It is guaranteed that there is a way to complete the task within $10^9$
days.

(In the original table the entries $20$, $10^5$ and the property A of
test cases $5 \sim 8$ are merged cells spanning several rows; they are
repeated on each row below.)

\begin{center}
\begin{tabular}{ccc}
\toprule
Test case number & $n \le$ & Special property \\
\midrule
$1$ & 20 & A \\
$2 \sim 4$ & 20 & none \\
$5 \sim 6$ & 500 & A \\
$7 \sim 8$ & $10^5$ & A \\
$9 \sim 10$ & $10^5$ & B \\
$11 \sim 13$ & $10^5$ & C \\
$14 \sim 16$ & $10^5$ & D \\
$17 \sim 20$ & $10^5$ & none \\
\bottomrule
\end{tabular}
\end{center}

Special property A: for all $1 \le i \le n$, $c_i = 0$.

Special property B: for all $1 \le i < n$, $u_i = i$ and $v_i = i + 1$.

Special property C: no plot has more than 2 roads directly connected to
it.

Special property D: for all $1 \le i < n$, $u_i = 1$.

\section*{Source}

The problems were translated from the official statement PDF
(\texttt{download.php?type=statement\&id=...}) displayed on the following
QOJ problem pages:
\begin{itemize}
\item Problem 1 (lock): \texttt{https://qoj.ac/problem/7813}
\item Problem 2 (game): \texttt{https://qoj.ac/problem/7814}
\item Problem 3 (struct): \texttt{https://qoj.ac/problem/7815}
\item Problem 4 (tree): \texttt{https://qoj.ac/problem/7816}
\end{itemize}
Contest page: \texttt{https://qoj.ac/contest/1428} (CSP-S 2023).

\end{document}
